\section[Stage 067]{Stage 067: Full Reduced Moving-Throat PDE Write-Up Skeleton and Remaining Theorem Gap}
\label{stage:067}

\stagefield{Part / anchor}{Part III; secondary anchor \resultanchor{MTDC-T7}.}
\claimstatus{\StatusExactClosure{} for the reduced hierarchy organization; \StatusOpen{} for final actual-branch realization.}
\stagefield{Purpose}{Stage~067 organizes the completed reduced hierarchy into a write-up-ready theorem chain and isolates the final outgoing quadrupole product as the remaining gap.}
\stagefield{Inputs}{Stages~020--066.}

\stagefield{Verification}{SymPy audit: none yet.  Mathematica audit: \StageFile{mathematica/moving_throat_pde_stage067_full_reduced_pde_writeup_mathematica_audit.wl}; transcript: \StageFile{mathematica/output/moving_throat_pde_stage067_full_reduced_pde_writeup_mathematica_audit.txt}.}

\subsection*{Derivation}

The reduced PDE chain is now
\begin{equation}
\label{eq:app-stage067-chain}
\begin{aligned}
\text{continuum placement}
&\rightarrow
\text{support completion}
\rightarrow
\text{transport/parent gain}
\\
&\rightarrow
\text{Family--1 support map}
\rightarrow
R_{\rm quad}.
\end{aligned}
\end{equation}
The remaining open datum is the outgoing branch's demand side, i.e. the conservative quadrupole module and its loading ratio.  The support/source machinery can be cited as completed within the declared reduced Family--1 closure once that ratio is supplied.

\stagefield{Output}{A structured reduced-PDE theorem chain and handoff to the outgoing-branch loading-ratio problem.}
\stagefield{Downstream use}{Stages~068--073 finish the support side by showing that the natural minimal module supplies a small ratio.}
